% !PRINTED 435
\weberpage{435}
% Unnumbered chapter-opening leaf; scan sequence fixes this leaf as i0466.

\begin{center}
\gesperrt{Eighteenth Chapter.}

\gesperrt{Galois Group of the Class Equation.}
\end{center}

\begin{center}\rule{0.18\linewidth}{0.4pt}\end{center}

\begin{center}
\textbf{\S\,120. \gesperrt{Relations between the class invariants of the same
discriminant.}}
\end{center}

For what follows it is convenient to have a notation by which the
dependence of the class invariant on the \gesperrt{class of forms} is expressed
more simply. We therefore set, when $k$ denotes an arbitrary
 primitive class of discriminant $-m$, which
according to § 118, III. is represented by $(a, b, c p^{\varepsilon})$:
\[
(k) = j\left(\frac{-b + \sqrt{-m}}{2a}\right) = j(\omega).
\]
When one or more of the $(k)$ occur under a function
 symbol, we shall also omit the brackets,
writing for example $f(k, k', \ldots)$ for $f[(k), (k'), \ldots]$.

If therefore
\weq{1}{(k) = j(\omega),}
and $l$ is a class represented by $(p, b, a c p^{\varepsilon-1})$, then by
§ 118, III.:
\weq{2}{(lk) = j\left(\frac{\omega}{p}\right),}
and so the two quantities
\[
u = (k), \quad v = (lk)
\]
satisfy the invariant equation
\weq{3}{F_p(v, u) = 0.}
The quantity $v = (lk)$ is, however, in addition a root of the
class equation
\weq{4}{H_m(v) = 0,}
and it is first to be established how many roots equation
(3) has in common with (4).
